\section{How we surveyed}
\label{sec:method}

\paragraph{Type of review.} This is a structured narrative review, not a
systematic review. We chose the families and the seed works from the
questions in \cref{sec:introduction}, then followed citations and searched for
recent work. We did not follow a registered protocol such as PRISMA, and we did
not compute inter-rater agreement.

\paragraph{Families.} We grouped approaches into seven families, by the problem
each was built to solve:
\begin{enumerate}[label=F\arabic*.]
  \item Responsible-AI practice: principles, design methods, documentation,
        audits, toolkits and runtime guardrails.
  \item Access control, policy engines and privacy compliance in data systems.
  \item Privacy-enhancing technologies for releasing data and statistics.
  \item Federated learning: algorithms, attacks, defences, frameworks and
        deployments.
  \item Bringing code to the data: federated analytics, distributed research
        networks, trusted research environments, data spaces, data clean rooms
        and privacy-preserving query compilers.
  \item Architectural precedents: moving computation to data, federated data
        architectures, separating decisions from effects, and declarative
        deployment.
  \item Release governance for machine learning, and the regulations and
        standards that set the requirements.
\end{enumerate}

\paragraph{Sources.} We searched the web and the publishers' indexes in
November 2026 and checked the authors, title, venue and year of every academic
work against Crossref or arXiv. For products and standards without a paper we
cite the vendor's or publisher's documentation and give the date we read it.
We include \SurveyWorks{} works in total.

\paragraph{Dimensions.} For each work we recorded seven dimensions, which refine
the four questions of \cref{sec:introduction}:
\begin{description}[font=\normalfont\bfseries, leftmargin=1.2em, style=nextline]
  \item[D1 What is decided] Permit or deny; a transformation (how much, in
        what form); a privacy mechanism and its parameters; who may take part;
        whether a result or model may be released.
  \item[D2 How the decision is expressed] Code in the application,
        configuration, a policy language, a person, a learned model, or a typed
        function; and whether it can be analysed or verified.
  \item[D3 Where it is enforced] A central store, each data source, the point
        of egress, or model release.
  \item[D4 Lifecycle stages covered] Reading data, computing or training,
        sending results, releasing a model, deploying a workload.
  \item[D5 Evidence] Audit logs, attestations, documentation, provenance.
  \item[D6 Failure behaviour] What happens when the decision point cannot
        answer: deny (fail closed) or allow (fail open).
  \item[D7 Coupling] To a domain, to application code, or to one vendor.
\end{description}

\paragraph{Threats to validity.} The review covers work in English that web
search and the publishers' indexes return. Industry products change often; we
describe each product as its documentation stood on the date we read it.
Because we built Dagents, we may read other work in its terms. To reduce this,
each family's section records what that family does that Dagents does not.
